\documentclass[11pt]{article}

\usepackage[margin=1in]{geometry}
\usepackage{amsmath,amssymb}
\usepackage{enumitem}
\usepackage{hyperref}
\usepackage{microtype}

\title{End-to-End Formalization: Compute-Efficient Structural Clipping for CISPO}
\author{}
\date{}

\begin{document}
\maketitle

\section{The Problem: The $O(N^2)$ Memory Wall in Structural RLHF}

Working backwards from empirical failures in reasoning tasks reveals that standard CISPO is artificially bottlenecked by \textbf{token-agnostic truncation}.

The standard objective applies a static scalar bound ($\epsilon_{\mathrm{high}}$) to the importance weight $r_t(\theta)$ for every token:
\begin{equation}
  L_{\mathrm{CISPO}}(\theta) = - \mathbb{E} \left[ \mathrm{detach}\!\left(\min(r_t(\theta), \epsilon_{\mathrm{high}})\right) \cdot \hat{A}_t \cdot \log \pi_\theta(a_t \mid s_t) \right]
\end{equation}

This mechanism clips the gradient of a critical reasoning token (e.g., a mathematical operator) exactly as it would a structural filler word, preventing the active policy ($\pi_\theta$) from rapidly updating critical logical hubs.

To solve this, the clipping bound must be dynamic ($\epsilon_t$), scaled by the token's structural importance (attention in-degree). However, calculating exact attention saliency requires caching the full $N \times N$ attention matrix across multiple layers during the active RL rollout. In a standard PPO/CISPO training loop, backpropagation already consumes maximum VRAM. Forcing $\pi_\theta$ to maintain an $O(N^2)$ spatial complexity for attention matrices triggers immediate Out-of-Memory (OOM) errors and destroys training throughput.

\section{The Solution: Decoupled \& Sub-Linear Saliency}

To resolve the bottleneck, we must extract structural saliency without caching $N \times N$ matrices on the active policy. We achieve this through a unified architecture that isolates computation and reduces spatial complexity.

\subsection{Component A: Reference Model Offloading (The Systems Shortcut)}

Instead of forcing the active policy ($\pi_\theta$) to compute its own saliency, we offload the calculation entirely to the frozen reference model ($\pi_{\mathrm{ref}}$). Because the foundational semantic routing of a prompt (the identification of logic gates and numbers) is structurally inherent to the sequence, the attention topology differs negligibly between $\pi_{\mathrm{ref}}$ and $\pi_\theta$. We extract all structural metrics exclusively during the $\pi_{\mathrm{ref}}$ forward pass (used for KL-divergence), adding zero FLOPs or memory overhead to the active policy.

\subsection{Component B: Sub-Linear Saliency Extraction}

Within the offloaded $\pi_{\mathrm{ref}}$ pass, we utilize one of two highly efficient metrics to quantify the raw structural importance score ($S_t$) for each token without materializing the full attention matrix.

\textbf{Path 1: The Architectural Proxy (Key-Vector $L_2$ Norms)}

Tokens that act as structural hubs naturally develop massive Key vectors so future queries can attend to them. We intercept the materialized KV-cache of $\pi_{\mathrm{ref}}$ at the final reasoning layers ($L_{\mathrm{top}}$) and compute the $L_2$ norm.
\begin{itemize}[leftmargin=*]
  \item \textbf{Spatial Complexity:} $O(1)$ overhead.
  \item \textbf{Metric:} $S_t = \|k_t^{(L_{\mathrm{top}})}\|_2$
\end{itemize}

\textbf{Path 2: The Exact Algorithmic Solution ($O(N)$ Streaming Accumulators)}

Instead of storing the full matrix, we dynamically accumulate the in-degree score. At each generation step $i$, $\pi_{\mathrm{ref}}$ computes a $1 \times i$ attention vector $\mathbf{a}_i$. We initialize a 1D tensor $\mathbf{c}$ of size $N$ and update the running total at each step: $\mathbf{c}_{1:i} \leftarrow \mathbf{c}_{1:i} + \mathbf{a}_i$.
\begin{itemize}[leftmargin=*]
  \item \textbf{Spatial Complexity:} $O(N)$ per head.
  \item \textbf{Metric:} $S_t = \mathbf{c}[t]$
\end{itemize}

\section{End-to-End Implementation Formalization}

The complete pipeline operates in three sequential steps during the post-training rollout.

\subsection{Step 1: Forward Pass \& Saliency Extraction}

During the generation and KL-divergence phase, the reference model $\pi_{\mathrm{ref}}$ processes the sequence. We extract the raw structural importance score $S_t$ for every token $t \in [1, T]$ using either the KV-norm proxy or the streaming accumulator.

\subsection{Step 2: Dynamic Bound Normalization}

To ensure the active policy's variance remains theoretically bounded over the entire sequence, we cannot simply use $S_t$ as the clipping threshold. We must normalize the raw scores to create a multiplier $m_t$, ensuring the mean of the dynamic bounds strictly equals the standard scalar $\epsilon_{\mathrm{high}}$.

We calculate the dynamic bound $\epsilon_t$ as follows:
\begin{align}
  m_t &= \frac{S_t}{\frac{1}{T} \sum_{j=1}^T S_j} \\
  \epsilon_t &= \epsilon_{\mathrm{high}} \cdot m_t
\end{align}

\textit{Implementation Note:} A lower bound $\delta$ can be applied ($m_t = \max(m_t, \delta)$) to prevent $\epsilon_t$ from collapsing entirely to zero on extreme filler tokens.

\subsection{Step 3: The AD-CISPO Objective Update}

With the frozen tensor $\epsilon_t$ successfully derived from $\pi_{\mathrm{ref}}$, it is passed to the active policy $\pi_\theta$. The standard scalar is replaced, yielding the final Adaptive and Dense CISPO objective:
\begin{equation}
  L_{\text{AD-CISPO}}(\theta) = - \mathbb{E} \left[ \mathrm{detach}\!\left(\min(r_t(\theta), \epsilon_t)\right) \cdot \hat{A}_t \cdot \log \pi_\theta(a_t \mid s_t) \right]
\end{equation}

\textbf{Result:} The active policy seamlessly executes gradient updates where critical tokens possess widened learning corridors, while filler tokens are aggressively clamped. The pipeline strictly maintains standard training throughput by eliminating all $O(N^2)$ tracking requirements.

\end{document}
